\documentclass[12pt,letterpaper]{article}
\usepackage[utf8]{inputenc}
\usepackage[spanish]{babel}
\usepackage{amsmath, amssymb, amsfonts}
\usepackage{graphicx}
\usepackage{hyperref}
\usepackage{enumitem}
\usepackage[margin=2.5cm]{geometry}

\title{Preparación para la sesión: Aritmética de punto flotante}
\author{Análisis Numérico}
\date{Sesión inversa}

\begin{document}

\maketitle

\begin{center}
\textbf{Objetivo:} Que al llegar a clase ya se hayan encontrado con el problema de que $0.2 + 0.3 \neq 0.5$ en una computadora.
\end{center}

\section*{1. Repaso de representación binaria}

En nuestra vida cotidiana usamos el sistema decimal (base 10). Sin embargo, las computadoras almacenan la información en binario (base 2). Para entender por qué ocurren errores al sumar números ``sencillos'', necesitamos recordar cómo se escriben los números fraccionarios en binario.

\subsection*{Números enteros en binario}
Un número entero se representa como una suma de potencias de 2:

\[
1011_2 = 1\cdot 2^3 + 0\cdot 2^2 + 1\cdot 2^1 + 1\cdot 2^0 = 8 + 0 + 2 + 1 = 11_{10}
\]

\subsection*{Números fraccionarios en binario}
Para la parte fraccionaria, usamos potencias negativas de 2:

\[
0.101_2 = 1\cdot 2^{-1} + 0\cdot 2^{-2} + 1\cdot 2^{-3} = \frac{1}{2} + 0 + \frac{1}{8} = 0.625_{10}
\]

\subsection*{Algoritmo para convertir una fracción decimal a binario}
Para convertir un número fraccionario a binario, multiplicamos sucesivamente por 2 y tomamos la parte entera:

\begin{enumerate}
    \item Multiplicar la fracción por 2.
    \item La parte entera del resultado es el primer dígito binario.
    \item Repetir con la parte fraccionaria restante.
\end{enumerate}

\subsection*{Ejemplos}

\subsubsection*{Ejemplo 1: Convertir $0.625_{10}$ a binario}
\begin{align*}
    0.625 \times 2 &= 1.25 && \text{parte entera } 1 \\
    0.25 \times 2  &= 0.5  && \text{parte entera } 0 \\
    0.5 \times 2   &= 1.0  && \text{parte entera } 1
\end{align*}
Por lo tanto,
\[
0.625_{10} = 0.101_2
\]

\subsubsection*{Ejemplo 2: Convertir $0.3_{10}$ a binario}
Este es el que nos interesa:
\begin{align*}
    0.3 \times 2 &= 0.6 && \text{parte entera } 0 \\
    0.6 \times 2 &= 1.2 && \text{parte entera } 1 \\
    0.2 \times 2 &= 0.4 && \text{parte entera } 0 \\
    0.4 \times 2 &= 0.8 && \text{parte entera } 0 \\
    0.8 \times 2 &= 1.6 && \text{parte entera } 1 \\
    0.6 \times 2 &= 1.2 && \text{parte entera } 1
\end{align*}
Observamos que el patrón se repite: $0.3_{10} = 0.01001100110011\ldots_2$, ¡es un número periódico en binario!

\subsubsection*{Ejemplo 3: Convertir $0.2_{10}$ a binario}
\begin{align*}
    0.2 \times 2 &= 0.4 && \text{parte entera } 0 \\
    0.4 \times 2 &= 0.8 && \text{parte entera } 0 \\
    0.8 \times 2 &= 1.6 && \text{parte entera } 1 \\
    0.6 \times 2 &= 1.2 && \text{parte entera } 1 \\
    0.2 \times 2 &= 0.4 && \text{parte entera } 0
\end{align*}
Nuevamente, el patrón se repite: $0.2_{10} = 0.001100110011\ldots_2$. 

\begin{center}
\fbox{\begin{minipage}{0.9\textwidth}
\textbf{Conclusión importante:} 
$0.3$ y $0.2$ son números con una representación binaria \textbf{infinita y periódica}. Como la computadora solo puede almacenar un número finito de dígitos, estos números se \textbf{redondean} o \textbf{truncan}. Esto introduce un pequeño error desde el momento en que se guardan.
\end{minipage}}
\end{center}

\section*{2. Errores absoluto y relativo}

\subsection*{Definiciones}

\begin{itemize}
    \item \textbf{Error absoluto:} 
    \[
    E_a = |\text{Valor real} - \text{Valor aproximado}|
    \]
    \item \textbf{Error relativo:}
    \[
    E_r = \frac{|\text{Valor real} - \text{Valor aproximado}|}{|\text{Valor real}|}
    \]
    (siempre que el valor real no sea cero)
\end{itemize}

El error relativo es más útil porque nos dice la magnitud del error en relación con el tamaño del número.

\subsection*{Ejemplo práctico:}
Si al aproximar $\pi = 3.14159265\ldots$ usamos $3.14$:
\begin{itemize}
    \item Error absoluto: $|3.14159265 - 3.14| = 0.00159265$
    \item Error relativo: $\frac{0.00159265}{3.14159265} \approx 0.000507$, es decir, un $0.05\%$
\end{itemize}

\section*{3. Tarea para antes de la sesión}

Para que lleguemos con material para discutir, realicen lo siguiente:

\begin{enumerate}
    \item Tomen una calculadora o abran \texttt{Julia} y calculen:
    \[
    0.2 + 0.3
    \]
    \item Comparen el resultado con $0.5$. En Julia, escriban:
    \begin{verbatim}
    0.2 + 0.3 == 0.5
    \end{verbatim}
    \item Reflexionen sobre el resultado. ¿Obtuvieron \texttt{true} o \texttt{false}?
    \item Investiguen: ¿a qué tipo de número en punto flotante corresponde el resultado de $0.2+0.3$? ¿Por qué no es exactamente $0.5$?
\end{enumerate}

\section*{4. Lo que haremos en clase}

En la sesión presencial:
\begin{itemize}
    \item Discutiremos sus hallazgos sobre $0.2 + 0.3$.
    \item Veremos cómo se almacena un número en el estándar IEEE 754 (punto flotante).
    \item Analizaremos ejemplos concretos de errores de redondeo en operaciones sencillas.
    \item Definiremos formalmente los errores absoluto y relativo.
    \item Haremos algunos ejercicios para que vean cómo estos errores afectan los métodos numéricos.
\end{itemize}

\section*{Pregunta final (para que piensen antes de llegar)}
¿Qué creen que pasaría si en una computadora sumamos $0.1$ diez veces? ¿Obtendremos exactamente $1.0$?

\vspace{1cm}
\textit{“El primer paso para dominar el análisis numérico es desconfiar de las respuestas que da la computadora.”}

\end{document}